% ===========================================================================
% Section 2: Related Work
% ===========================================================================

% ---------------------------------------------------------------------------
\subsection{Retrieval-Augmented Generation}
\label{subsec:rw_rag}
% ---------------------------------------------------------------------------

Retrieval-Augmented Generation~\cite{lewis2020rag} grounds LLM outputs in externally retrieved evidence, combining a retriever with a generator.
Early approaches use a single retrieval pass with dense retrievers such as DPR~\cite{karpukhin2020dense} and ColBERT~\cite{khattab2020colbert}, both building on contextualized BERT~\cite{devlin2019bert} representations, or sparse methods like BM25~\cite{robertson2009bm25}.
Pre-training with retrieval objectives (REALM~\cite{guu2020realm}) and retrieval-augmented generation architectures (Fusion-in-Decoder~\cite{izacard2021leveraging}, Atlas~\cite{izacard2023atlas}) further improve knowledge grounding.
Recent work improves retrieval quality through query rewriting~\cite{ma2023query}, hypothetical document embeddings (HyDE)~\cite{gao2023hyde}, LLM-augmented dense retrieval~\cite{peng2025sptar,silva2024improving}, and hybrid retrieval combining dense and sparse methods via Reciprocal Rank Fusion (RRF)~\cite{cormack2009rrf}.
For a comprehensive survey of RAG methods, we refer readers to Gao et al.~\cite{gao2024ragsurvey}.

Our work builds on this foundation by employing hybrid FAISS~\cite{johnson2019faiss} + BM25 retrieval with RRF fusion as the base retrieval mechanism, enhanced by agentic refinement that iteratively improves retrieval quality through tool-augmented reasoning.

% ---------------------------------------------------------------------------
\subsection{Self-Corrective RAG}
\label{subsec:rw_corrective}
% ---------------------------------------------------------------------------

Self-corrective approaches introduce feedback loops that evaluate and refine retrieved passages.

\textbf{Self-RAG}~\cite{asai2024selfrag} trains a single LM to generate special reflection tokens that trigger on-demand retrieval and self-assess output quality.
While effective, it requires fine-tuning a custom model with reflection tokens, limiting flexibility.

\textbf{CRAG}~\cite{yan2024crag} introduces a lightweight retrieval evaluator that classifies documents as correct, incorrect, or ambiguous, triggering corrective actions including web search fallback and knowledge refinement through strip decomposition.
However, CRAG's single-pass correction with binary evaluation provides limited refinement granularity.

\textbf{Adaptive-RAG}~\cite{jeong2024adaptiverag} routes questions to different retrieval strategies based on estimated complexity, using a classifier trained to distinguish simple, moderate, and complex questions.
This complexity-aware approach is complementary to our work; our findings similarly show complexity-dependent improvement patterns.

Multi-hop dense retrieval methods offer an alternative approach: MDR~\cite{xiong2021mdr} performs iterative retrieval by encoding the question with previously retrieved passages, Baleen~\cite{khattab2021baleen} uses condensed summarization between retrieval hops, and BeamDR~\cite{zhao2021beamdr} applies beam search in the dense representation space.
Datasets such as IIRC~\cite{ferguson2020iirc} further test retrieval under incomplete information.

Other iterative retrieval approaches include \textbf{FLARE}~\cite{jiang2023flare}, which triggers retrieval when the model's generation confidence drops below a threshold;
\textbf{IRCoT}~\cite{trivedi2023ircot}, which interleaves retrieval with chain-of-thought reasoning steps for multi-hop questions;
and \textbf{Self-Ask}~\cite{press2023selfask}, which decomposes questions into sub-questions answered sequentially.
Verify-and-Edit~\cite{zhao2023verify} post-hoc verifies chain-of-thought reasoning by retrieving evidence for each reasoning step, while REPLUG~\cite{shi2024replug} treats the LM as a black box and prepends retrieved documents to the input.

Our approach differs from these methods in three aspects: (1)~we use an autonomous \emph{agent} rather than fixed correction rules; (2)~the agent has access to \emph{multiple specialized tools} beyond simple re-retrieval; and (3)~the entire pipeline is implemented as a \emph{declarative program} amenable to automatic optimization.

Table~\ref{tab:related_comparison} summarizes the key differences.

\begin{table*}[t]
\centering
\caption{Comparison of self-corrective RAG approaches. \ours{} uniquely combines multi-tool agentic refinement with declarative pipeline optimization.}
\label{tab:related_comparison}
\resizebox{\textwidth}{!}{%
\begin{tabular}{lccccc}
\toprule
\textbf{Method} & \textbf{Iterative} & \textbf{Tool Use} & \textbf{Declarative} & \textbf{Multi-hop} & \textbf{Agent} \\
\midrule
Naive RAG         & --          & --               & --       & --       & --        \\
Self-RAG~\cite{asai2024selfrag} & Selective   & --               & --       & \checkmark & Reflection \\
CRAG~\cite{yan2024crag}         & Single-pass & Web search       & --       & --       & Rules     \\
FLARE~\cite{jiang2023flare}     & Adaptive    & --               & --       & --       & Confidence \\
IRCoT~\cite{trivedi2023ircot}   & Interleaved & --               & --       & \checkmark & CoT       \\
Adaptive-RAG~\cite{jeong2024adaptiverag} & Selective & --          & --       & \checkmark & Router    \\
\textbf{\ours{}}  & \textbf{Adaptive} & \textbf{6 tools} & \textbf{DSPy}  & \checkmark & \textbf{ReAct}    \\
\bottomrule
\end{tabular}
}% end resizebox
\end{table*}

% ---------------------------------------------------------------------------
\subsection{Tool-Augmented Language Model Agents}
\label{subsec:rw_agents}
% ---------------------------------------------------------------------------

Tool-augmented LLMs extend language model capabilities through external tool invocation.
MRKL~\cite{karpas2022mrkl} proposes a modular neuro-symbolic architecture combining LLMs with external knowledge sources and discrete reasoning modules.
Toolformer~\cite{schick2023toolformer} trains models to autonomously decide when and how to call tools via self-supervised learning.
ART~\cite{paranjape2023art} automatically generates multi-step reasoning programs that interleave reasoning with tool calls.
ReAct~\cite{yao2023react} interleaves reasoning traces with actions in a Thought--Action--Observation loop, enabling models to plan, execute, and refine their approach dynamically.
Related reasoning paradigms include Tree of Thoughts~\cite{yao2023tree}, which explores multiple reasoning paths via tree search, and RAP~\cite{hao2023reasoning}, which frames LM reasoning as planning with a world model.
Reflexion~\cite{shinn2023reflexion} extends this with verbal reinforcement learning, where agents reflect on failures to improve subsequent attempts.

Broader agent frameworks have also emerged: HuggingGPT~\cite{shen2023hugginggpt} orchestrates specialized AI models as tools, ToolLLM~\cite{qin2024toolllm} facilitates tool mastery across 16,000+ APIs, and AgentBench~\cite{liu2024agentbench} provides a benchmark for evaluating LLMs as autonomous agents across diverse environments.
Additional benchmarks such as SWE-bench~\cite{jimenez2024swebench} and WebArena~\cite{zhou2024webarena} evaluate agents in software engineering and web navigation contexts, respectively.

In the RAG context, agentic approaches have emerged that use LLM agents to orchestrate retrieval~\cite{singh2025agenticrag}.
Recent work explores multi-step retrieval agents that decompose complex queries~\cite{khot2023decomposed}, perform iterative search, and synthesize evidence from multiple sources.
However, most existing agentic RAG systems focus on a single tool (search) with fixed orchestration logic, rather than providing a diverse tool set with autonomous tool selection.

\ours{} equips the ReAct agent with six specialized tools, enabling richer retrieval strategies that go beyond simple re-querying: query decomposition for multi-hop questions, quality evaluation for informed refinement, progressive passage inspection, and domain-specific tools for enterprise documents.

% ---------------------------------------------------------------------------
\subsection{Declarative Language Model Programming}
\label{subsec:rw_dspy}
% ---------------------------------------------------------------------------

The broader field of prompt optimization has explored various approaches to automatically improve LM prompts.
AutoPrompt~\cite{shin2020autoprompt} uses gradient-guided search over discrete tokens, APE~\cite{zhou2023ape} treats prompt generation as a natural language synthesis problem, and OPRO~\cite{yang2024opro} uses LLMs themselves as optimizers.
More recently, TextGrad~\cite{yuksekgonul2025textgrad} backpropagates natural language feedback through computational graphs for end-to-end optimization.

DSPy~\cite{khattab2024dspy}, building on the Demonstrate-Search-Predict framework~\cite{khattab2023dsp}, introduces a programming model for LM-based systems based on declarative \emph{Signatures} (typed input/output specifications) and composable \emph{Modules} (\eg, ChainOfThought, ReAct).
This separates the program logic from the prompt implementation, enabling automatic optimization through compilers like BootstrapFewShot and MIPROv2~\cite{opsahl2024miprov2}.

BootstrapFewShot collects successful execution traces and uses them as few-shot demonstrations, automatically generating in-context examples.
MIPROv2 extends this by jointly optimizing natural language instructions and demonstrations using Bayesian surrogate models.

While DSPy has been applied to various NLP tasks, its application to self-corrective RAG pipelines---particularly in combination with ReAct agents and multi-tool orchestration---has not been systematically studied.
Our work demonstrates that DSPy's structured prompting provides substantial benefits even without optimization ($+0.1$--$0.2$ F1), with additional gains from automatic optimization that are dataset-dependent.
